\chapter{Problem Analysis and Dataset}
\label{chap:problem-dataset}

This chapter answers two questions and keeps them apart. How large is the
reproducibility problem in general, and what exactly does this thesis's
implementation operate on? The first is answered by numbers from the
original, full-scale FAIR Jupyter study. The second is answered by a much
smaller, precisely defined dataset drawn from the newer Docker-based
pipeline. The two are related, because the second is a residual subset of
the same kind of failure, but they are not the same numbers and are never
mixed here.

\section{Scale of the Problem}
\label{sec:pd-l1}

The original FAIR Jupyter study executed every notebook it collected using a
fresh conda environment per repository~\cite{samuel2024gigascience}. The
figures below were recomputed for this thesis from the published database
of that study~\cite{samuel2024zenodo}, using the local snapshot of its
\texttt{executions} table. A notebook counts as executed when its
execution record is neither an environment-installation failure nor a
skipped notebook, and the two exception counts are taken from the recorded
failure reason of each execution:

\begin{itemize}
  \item Total notebooks executed: \textbf{10,389}
  \item \texttt{ModuleNotFoundError}: 5,562 (53.5\%)
  \item \texttt{ImportError}: 1,014 (9.8\%)
  \item Combined dependency-related failures: \textbf{6,576 (63.3\%)}
\end{itemize}

The published article itself reports 10,388 notebooks attempted for
re-execution and 6,588 \texttt{ModuleNotFoundError} and
\texttt{ImportError} exceptions, computed over a different
denominator~\cite{samuel2024gigascience}; the small differences reflect
the counting rule stated above, not a different dataset. In other words,
roughly two out of every three executed notebooks failed for a
dependency-related reason. The missing modules were dominated by
biomedical and machine-learning packages, consistent with the PubMed Central
origin of the dataset: \texttt{anndata} (1258 occurrences),
\texttt{scanpy} (423), \texttt{pandas} (177), \texttt{cemba\_data} (176),
\texttt{tensorflow} (149), \texttt{Bio} (108), and \texttt{fastai} (103) were
the most frequently missing modules in that snapshot.

These numbers exist to establish that dependency failure is a large-scale,
systematic problem, not an edge case worth ignoring. \textbf{They are not the
dataset this thesis's implementation runs on.} They come from the earlier
conda-based run of the pipeline, before the Docker-based extension described
next.

\section{The Working Dataset}
\label{sec:pd-i1}

\subsection{Why the Docker Pipeline, Not the Conda Run}
\label{sec:pd-why-docker}

The newer, containerised pipeline~\cite{samuel2025docker} re-executes
each repository inside its own Docker container. Containerisation
resolves many of the failures counted in \Cref{sec:pd-l1} that were
really caused by a missing system tool or an operating-system mismatch,
not by the notebook's Python dependencies. What is left over after
containerisation -- notebooks that still fail even inside a container
built specifically for them -- is a cleaner, more targeted signal of a
genuine dependency problem. This is why the Docker pipeline, not the
original conda run, was chosen as the integration target. Its residual
failures are exactly the gap this thesis is meant to close. And its
result schema already tags each failure with a pre-computed
\texttt{error\_category}, which removes the need to re-implement
failure detection from scratch.

\subsection{Source and Filtering}
\label{sec:pd-source}

The working dataset is derived from the Docker pipeline's own result
database (\texttt{notebook\_executions} table). Its existing categorisation
routine already maps two exception types into a bucket called
\texttt{DEPENDENCY\_ERROR}: \texttt{ModuleNotFoundError} and
\texttt{ImportError}. These two types were kept, unchanged, as this thesis's
scope, for two reasons: they are exactly what the upstream pipeline already
flags as dependency-related, and, taken together, they cover both a package
being completely absent and a package being present but missing a specific
name the notebook tries to import.

Filtering every row where \texttt{error\_category = 'DEPENDENCY\_ERROR'}
yields \textbf{214} rows: 172 \texttt{ModuleNotFoundError} and 42
\texttt{ImportError}. This extraction and every count in this section is
produced by \texttt{scripts/prepare\_dependency\_dataset.py}, and is written
out as a CSV file plus a machine-readable statistics file so that the
numbers can be independently recomputed.

\subsection{Repair Scope: Why pip-Only}
\label{sec:pd-scope}

Not every one of the 214 rows can be repaired the same way. Version~1 of the
repair layer is restricted, deliberately, to dependency problems that a
single \texttt{pip install} command can fix: installing a missing package,
or pinning a specific version. This scope decision was made for a concrete
reason: it keeps the first version of the system controlled and measurable.
A repair action that also has to reason about system-level packages
(installed with \texttt{apt}, not \texttt{pip}) or about renaming an import
path inside the notebook's own code introduces a second, much larger kind of
decision that this thesis does not attempt to automate yet.

Applying this scope splits the 214 rows into the four subtypes of
\Cref{tab:pd-subtypes}, and those in turn into two groups.

\begin{table}[htbp]
  \centering
  \begin{tabular}{|l|r|p{6.3cm}|}
    \hline
    \textbf{Subtype} & \textbf{Count} & \textbf{Meaning} \\
    \hline
    \texttt{missing\_package} & 179 & The imported module is genuinely absent. \\
    \hline
    \texttt{wrong\_version} & 21 & The module is present, but the notebook imports a name that no longer exists in the installed version. \\
    \hline
    \texttt{system\_library} & 10 & A missing shared object file (e.g.\ \texttt{libxcb.so.1}), installable only with \texttt{apt}, not \texttt{pip}. \\
    \hline
    \texttt{mapping\_unknown} & 4 & An ambiguous local or path-shadowed name (e.g.\ \texttt{utils}) that cannot safely be mapped to a PyPI package. \\
    \hline
  \end{tabular}
  \caption{Dependency-error subtypes across the 214-row working dataset.}
  \label{tab:pd-subtypes}
\end{table}

The 10 \texttt{system\_library} rows are excluded because they need
\texttt{apt}-level installation, which is outside the pip-only scope by
definition. A missing shared library cannot be fixed by any pip
command, however it is phrased. The 4 \texttt{mapping\_unknown} rows are
excluded because their failing name (for example a generic module
called \texttt{utils}) cannot be safely resolved to one specific PyPI
package without guessing. Guessing a package name is exactly the kind
of unfounded suggestion this thesis's PyPI-grounding design
(\Cref{chap:architecture}) is built to avoid. This leaves \textbf{200
usable} rows and \textbf{14 excluded} rows.

\begin{table}[htbp]
  \centering
  \begin{tabular}{|l|r|p{7.5cm}|}
    \hline
    \textbf{Split} & \textbf{Count} & \textbf{Purpose} \\
    \hline
    \texttt{dev} & 13 & Used to design and manually sanity-check the explanation and repair prompts (\Cref{chap:methodology}). \\
    \hline
    \texttt{evaluation} & 187 & Usable rows reserved for the final repair evaluation, reported in \Cref{chap:validation}. \\
    \hline
    \texttt{excluded} & 14 & Within explanation scope; excluded from automated repair (\Cref{sec:pd-explanation-scope}). \\
    \hline
  \end{tabular}
  \caption{Split policy over the 214-row working dataset.}
  \label{tab:pd-splits}
\end{table}

\Cref{tab:pd-splits} shows how the 214 rows divide across the three
resulting splits. The 13-row \texttt{dev} split was used exclusively during
implementation, to design and manually sanity-check the explanation and
repair prompts and to run the small real-system pilots reported in
\Cref{chap:validation}. The
187-row \texttt{evaluation} split was reserved from prompt and
repair-strategy tuning: no prompt, mapping entry, or repair rule was
designed against it. The results reported in \Cref{chap:validation} were
generated using the final frozen implementation described in
\Cref{chap:methodology}.

\subsection{Explanation Scope Is Wider Than Repair Scope}
\label{sec:pd-explanation-scope}

One further decision affects every later chapter and is worth stating
clearly. The pip-only restriction above governs \emph{repair} eligibility
only. It does not restrict \emph{explanation}. All 214 rows, including the
14 excluded ones, stay within the explanation scope defined for objective
O1. Only the 200 usable rows are additionally in scope for the
repair-proposal and fix-application components (objectives O2 and O3). A
researcher whose notebook fails because of a missing system library still
deserves to be told what went wrong, even though the current repair layer
cannot fix that case automatically.

\Cref{fig:pd-scope} shows this scope boundary in one place. Every one of
the 214 records enters classification and stays within the explanation
objective, while only 200 are repair-eligible and feed the rest of the
repair layer.

\begin{figure}[h]
  \centering
  \resizebox{0.92\textwidth}{!}{%
  \begin{tikzpicture}[
      font=\small,
      comp/.style={draw, rounded corners, fill=blue!6, minimum width=3.1cm,
                   minimum height=1.05cm, align=center, node distance=8mm},
      elig/.style={draw, rounded corners, fill=green!8, minimum width=2.9cm,
                   minimum height=1.2cm, align=center, node distance=6mm},
      excl/.style={draw, rounded corners, fill=red!6, minimum width=2.9cm,
                   minimum height=1.2cm, align=center, node distance=6mm},
      leaf/.style={draw, rounded corners, fill=gray!10, minimum width=2.9cm,
                   minimum height=0.9cm, align=center, node distance=6mm},
      arr/.style={-{Latex[length=2.2mm]}, thick},
    ]
    % ---------- row of four subtypes ----------
    \node[elig] (c1) {\texttt{missing\_package}\\179\\\footnotesize repair-eligible};
    \node[elig, right=of c1] (c2) {\texttt{wrong\_version}\\21\\\footnotesize repair-eligible};
    \node[excl, right=26mm of c2] (c3) {\texttt{system\_library}\\10\\\footnotesize excluded from repair};
    \node[excl, right=of c3] (c4) {\texttt{mapping\_unknown}\\4\\\footnotesize excluded from repair};

    \coordinate (topmid) at ($(c1)!0.5!(c4)$);
    \node[comp, above=30mm of topmid] (records)
      {214 \texttt{DEPENDENCY\_ERROR} records\\(172 \texttt{ModuleNotFoundError} + 42 \texttt{ImportError})};
    \node[font=\small, below=8mm of records, text width=6cm, align=center] (classify)
      {Error classification / repair eligibility};

    \coordinate (busY) at ($(c1.north)+(0,7mm)$);
    \draw[arr] (records) -- (classify);
    \draw (classify.south) -- (classify.south |- busY);
    \draw (c1.north -| busY) -- (c4.north -| busY);
    \draw[arr] (c1.north -| busY) -- (c1.north);
    \draw[arr] (c2.north -| busY) -- (c2.north);
    \draw[arr] (c3.north -| busY) -- (c3.north);
    \draw[arr] (c4.north -| busY) -- (c4.north);

    % ---------- aggregate row ----------
    \coordinate (usablemid) at ($(c1)!0.5!(c2)$);
    \coordinate (exclmid) at ($(c3)!0.5!(c4)$);
    \node[comp, below=16mm of usablemid] (usable) {200 repair-usable};
    \node[comp, below=16mm of exclmid] (excluded) {14 excluded from repair};

    \draw[arr] (c1.south) -- (usable.north);
    \draw[arr] (c2.south) -- (usable.north);
    \draw[arr] (c3.south) -- (excluded.north);
    \draw[arr] (c4.south) -- (excluded.north);

    % ---------- leaves ----------
    \node[leaf, below left=12mm and 2mm of usable] (dev) {13 development};
    \node[leaf, below right=12mm and 2mm of usable] (evalset) {187 frozen evaluation};
    \draw[arr] (usable) -- (dev);
    \draw[arr] (usable) -- (evalset);

    \node[leaf, below left=12mm and 2mm of excluded] (sys) {10 \texttt{system\_library}};
    \node[leaf, below right=12mm and 2mm of excluded] (map) {4 \texttt{mapping\_unknown}};
    \draw[arr] (excluded) -- (sys);
    \draw[arr] (excluded) -- (map);
  \end{tikzpicture}%
  }
  \caption{Scope boundary over the 214-row working dataset. Every record
  enters classification and stays within the explanation objective O1, but
  only \texttt{missing\_package} and \texttt{wrong\_version} records are
  repair-eligible. Those 200 records split into the development and frozen
  evaluation sets.}
  \label{fig:pd-scope}
\end{figure}

\subsection{The Most Common Failing Modules}
\label{sec:pd-top-modules}

\begin{table}[htbp]
  \centering
  \begin{tabular}{|l|r||l|r|}
    \hline
    \textbf{Module} & \textbf{Occurrences} & \textbf{Module} & \textbf{Occurrences} \\
    \hline
    \texttt{pkg\_resources} & 22 & \texttt{google} & 5 \\
    \hline
    \texttt{sklearn} & 16 & \texttt{gnomad} & 5 \\
    \hline
    \texttt{rpy2} & 15 & \texttt{omero} & 5 \\
    \hline
    \texttt{numpy} & 14 & \texttt{dms\_variants} & 4 \\
    \hline
    \texttt{cana} & 12 & \texttt{bindingcalculator} & 4 \\
    \hline
    \texttt{celloracle} & 12 & \texttt{libXrender.so.1} & 4 \\
    \hline
    \texttt{keras} & 7 & \texttt{statsmodels} & 4 \\
    \hline
    \texttt{scipy} & 7 & \texttt{openslide-bin} & 4 \\
    \hline
    \texttt{skimage} & 6 & \texttt{libxcb.so.1} & 3 \\
    \hline
    \texttt{physicool} & 5 & \texttt{bokeh} & 3 \\
    \hline
  \end{tabular}
  \caption{The 20 most frequent failing modules in the 214-row working
  dataset, in descending order down the left column and continuing down the
  right. Two entries (\texttt{libXrender.so.1}, \texttt{libxcb.so.1}) are
  shared system libraries rather than Python modules, and fall outside the
  pip-only repair scope.}
  \label{tab:pd-top-modules}
\end{table}

Several of these entries already hint at a recurring implementation problem
addressed in \Cref{chap:methodology}: \texttt{sklearn} is not the PyPI
package name (it is \texttt{scikit-learn}), and \texttt{pkg\_resources} is
not an installable package at all but a module bundled with
\texttt{setuptools}. A repair layer that assumed a Python import name is
always also its own pip package name would fail on exactly the most common
cases in this dataset.

\section{Summary}
\label{sec:pd-summary}

The reproducibility problem is large at the scale of the full conda-based
FAIR Jupyter run (\Cref{sec:pd-l1}: 6,576 dependency-related failures out of
10,389 executed notebooks in the local snapshot of the published
dataset). This thesis does not attempt to repair that
entire set. It operates on a precisely defined, much smaller residual
subset: the 214 dependency-error rows still present after the newer Docker
pipeline's containerisation, of which 200 fall inside the pip-only v1 repair
scope. Of those 200 rows, 13 were used only for prompt design and small
real-system pilots, and the remaining 187 were reserved for the final
repair evaluation reported in \Cref{chap:validation}.

Keeping these numbers
distinct throughout the rest of this thesis avoids overstating what has
actually been measured.
